%
%%%%%%%%%%%%%%%% General document %%%%%%%%%%%%%%%%
\usepackage{graphicx}
\usepackage{hyperref}
\usepackage[latin1]{inputenc}
\usepackage{color}
\usepackage[french]{babel}
\usepackage{amsmath, amsfonts}
\usepackage{textgreek}
%\usepackage{siunitx}
\usepackage{rotating}
\usepackage{booktabs} % Allows the use of \toprule, \midrule and \bottomrule in tables
%\usepackage{natbib}
%\usepackage{latexsym}
%\usepackage{amssymb}
%
\usepackage{bm}  % bold math
\usepackage{booktabs}  % better tables
\usepackage{multirow}  % combine table rows
\usepackage{makecell}  % arrange content in table cells incl. linebreaks
\usepackage{csquotes}  % quotation/citation commands

%%%%%%%%%%%%%%%% FIGURES %%%%%%%%%%%%%%%%
\usepackage{subfigure}
\usepackage{epsfig} 
\usepackage{epstopdf}
%\usepackage{movie15}
%\usepackage{eurosym}
%\usepackage{gensymb}
%\usepackage{lmodern}
%\usepackage{tabularx}
\usepackage{booktabs}
\usepackage{xcolor}
\usepackage{colortbl}
% Define colors
\definecolor{rowlight}{gray}{0.92}
\definecolor{rowdark}{gray}{0.80}
%\usepackage{float}
%
%%%%%%%%%%%%%%% TIKZ %%%%%%%%%%%%%%%%%%%%
\usepackage{smartdiagram}
\usepackage{metalogo}
%\usepackage{dtklogos}
%
%%%%%%%%%%%%%%%% COLORS %%%%%%%%%%%%%%%%
\definecolor{violet}{rgb}{0.3,0.2,0.6}
\definecolor{violetpale}{rgb}{0.7,0.7,0.8}
\definecolor{marine}{RGB}{25,25,112}
\definecolor{marinepale}{rgb}{0.4,0.4,0.5}
\definecolor{myblue}{rgb}{0.2,0.2,0.7}
\definecolor{mycream}{rgb}{0.9712 0.940 0.807}
%
%
%%%%%%%%%%%%%%%% TEXT BOXES %%%%%%%%%%%%%%%%
\usepackage[absolute,showboxes,overlay]{textpos}
%\TPshowboxestrue % commenter une fois fini
\TPshowboxesfalse % d�commenter pour faire disparaitre les boites
\textblockorigin{10mm}{10mm} % origine des positions
%
%
%%%%%%%%%%%%%%%% BEAMER PACKAGE AND OPTIONS %%%%%%%%%%%%%%%%
%%% Beamer Layout %%%%%%%%%%%%%%%%%%%%%%%%%%%%%%%%%%
%%\usetheme{Madrid}
%%\usetheme{Warsaw}
\usetheme{Frankfurt}
%%\usetheme{Amsterdam}
\setbeamerfont{caption}{size=\scriptsize}
\setbeamertemplate{caption}[numbered]
%%\useoutertheme[subsection=false,shadow]{miniframes}
%%\setbeamertemplate{mini frames}{}
\useinnertheme{default}
\usefonttheme{serif}
\usepackage{palatino}
%%\setbeamertemplate{navigation symbols}{}
\setbeamerfont{frametitle}{size=\small}
\setbeamertemplate{footline}[page number]
%\setbeamercolor*{normal text}{fg=marine,bg=mycream} 
\setbeamercolor*{normal text}{fg=marine,bg=white} 
\setbeamercolor*{alerted text}{fg=red} 
\setbeamercolor*{example text}{fg=black} 
\setbeamercolor*{structure}{fg=marine}  
\setbeamercolor*{palette tertiary}{fg=black,bg=black!10} 
\setbeamercolor*{palette quaternary}{fg=black,bg=black!10} 
\setbeamertemplate{caption}{\raggedright\insertcaption\par}
%%%%%%%%%%%%%%%%%%%%%%%%%%%%%%%%%%%%%%%%%%%%%%%%%%%
\makeatletter
    \newenvironment{withoutheadline}{
        \setbeamertemplate{headline}[default]
        \def\beamer@entrycode{\vspace*{-\headheight}}
    }{}
\setbeamertemplate{mini frames}{}
\setbeamertemplate{headline}{}
\setbeamertemplate{navigation symbols}{}
\definecolor{myred}{rgb}{0.51,0.0,0.0} 
\setbeamercolor{frametitle}{fg=white, bg=myred}
\setbeamercolor{title}{fg=white, bg=myred}
%
%\newlength\leftsidebar
%\newlength\rightsidebar
%\makeatletter
%\setlength\leftsidebar{\beamer@leftsidebar}
%\setlength\rightsidebar{\beamer@rightsidebar}
%%\setbeamercovered{transparent}
%\def\beamertemplatetransparentcoveredmedium{\setbeamercovered{transparent=0.9}}
%\beamertemplatetransparentcoveredmedium
%%
%%\setbeamertemplate{headline}{%
%%\leavevmode%
%%  \hbox{%
%%    \begin{beamercolorbox}[wd=\paperwidth,ht=2.5ex,dp=1.125ex]{palette quaternary}%
%%    \insertsectionnavigationhorizontal{\paperwidth}{}{\hskip0pt plus1filll}
%%    \end{beamercolorbox}%
%%  }
%%}
%%\setbeamertemplate{headline}
%%{%
%%  \leavevmode%
%%  \begin{beamercolorbox}[wd=.3\paperwidth,ht=2.5ex,dp=1.125ex]{section in head/foot}%
%%    \hbox to .3\paperwidth{\hfil\insertsectionhead\hfil}
%%  \end{beamercolorbox}%
%%  \begin{beamercolorbox}[wd=.3\paperwidth,ht=2.5ex,dp=1.125ex]{subsection in head/foot}%
%%    \hbox to .3\paperwidth{\hfil\insertsubsectionhead\hfil}
%%  \end{beamercolorbox}%
%%}
%%\setbeamertemplate{navigation symbols}{} 
%%
%%\setbeamercolor*{lower separation line head}{bg=DeepSkyBlue4} 
%
%
%
%%%%%%%%%%%%%%%% TIKZ %%%%%%%%%%%%%%%%
\usepackage{tikz}
\usetikzlibrary{tikzmark}
\usetikzlibrary{decorations.pathreplacing,calc}
\usetikzlibrary{decorations.fractals}
\usepackage{pgfplots}
\usepackage{pgfplotstable}
%
%



%%%%%%%%%%%%%%%%%%%%%%%%%%%%%%%%%%%%%%%%%%%%%%%%%%%%%%%
%%
%
%\newcommand\MyText[4][0pt]{%
%\begin{tikzpicture}[overlay, remember picture]
%\draw [decoration={brace},decorate,thick]
%   ( $ ({pic cs:#3}|-{pic cs:#2}) + (#1,1.3ex) $ ) --  
%    node[anchor=west,xshift=5pt,text width=5cm] {#4} 
%  ( $ (pic cs:#3) + (#1,0) $ );
%\end{tikzpicture}%
%}

%
%\makeatother
%
%% See if the author has AMS Euler fonts installed: If they have, attempt
%% to use the 'upmath' package to provide upright math.
%%%\usepackage{lipsum}
%%\usepackage{subfigure}
%%\linespread{1.6}
%%\textwidth 14.7 true cm
%%\textheight 21.6 true cm
%
%%\ifCUPmtlplainloaded \else
%%  \checkfont{eurm10}
%%  \iffontfound
%%    \IfFileExists{upmath.sty}
%%      {\typeout{^^JFound AMS Euler Roman fonts on the system,
%%                   using the 'upmath' package.^^J}%
%%       \usepackage{upmath}}
%%      {\typeout{^^JFound AMS Euler Roman fonts on the system, but you
%%                   dont seem to have the}%
%%       \typeout{'upmath' package installed. JFM.cls can take advantage
%%                 of these fonts,^^Jif you use 'upmath' package.^^J}%
%%       \providecommand\upi{\pi}%
%%      }
%%  \else
%%    \providecommand\upi{\pi}%
%%  \fi
%%\fi 
%
%%% See if the author has AMS symbol fonts installed: If they have, attempt
%%% to use the 'amssymb' package to provide the AMS symbol characters.
%
%%\ifCUPmtlplainloaded \else
%%  \checkfont{msam10}
%%  \iffontfound
%%    \IfFileExists{amssymb.sty}
%%      {\typeout{^^JFound AMS Symbol fonts on the system, using the
%%                'amssymb' package.^^J}%
%%       \usepackage{amssymb}%
%%       \let\le=\leqslant  \let\leq=\leqslant
%%       \let\ge=\geqslant  \let\geq=\geqslant
%%      }{}
%%  \fi
%%\fi
%
%%% See if the author has the AMS 'amsbsy' package installed: If they have,
%%% use it to provide better bold math support (with \boldsymbol).
%
%%\ifCUPmtlplainloaded \else
%%  \IfFileExists{amsbsy.sty}
%%    {\typeout{^^JFound the 'amsbsy' package on the system, using it.^^J}%
%%     \usepackage{amsbsy}}
%%    {\providecommand\boldsymbol[1]{\mbox{\boldmath $##1$}}}
%%\fi
%
%
%%%% Example macros (some are not used in this sample file) %%%
%
%%Comments
%\def\NOTE#1{{\textcolor{red}{\bf [#1]}}}  % note
%\def\ADD#1{{\textcolor{blue}{\bf #1}}}    % question
%
%% For units of measure
%\newcommand\dynpercm{\nobreak\mbox{$\;$dyn\,cm$^{-1}$}}
%\newcommand\cmpermin{\nobreak\mbox{$\;$cm\,min$^{-1}$}}
%
%% Various bold symbols
%\providecommand\bnabla{\boldsymbol{\nabla}}
%\providecommand\bcdot{\boldsymbol{\cdot}}
%\newcommand\biS{\boldsymbol{S}}
%\newcommand\etb{\boldsymbol{\eta}}
%
%% For multiletter symbols
%\newcommand\Real{\mbox{Re}} % cf plain TeX's \Re and Reynolds number
%\newcommand\Imag{\mbox{Im}} % cf plain TeX's \Im
%\newcommand\Rey{\mbox{\textit{Re}}}  % Reynolds number
%\newcommand\Pran{\mbox{\textit{Pr}}} % Prandtl number, cf TeX's \Pr product
%\newcommand\Pen{\mbox{\textit{Pe}}}  % Peclet number
%\newcommand\Ai{\mbox{Ai}}            % Airy function
%\newcommand\Bi{\mbox{Bi}}            % Airy function
%
%% For sans serif characters:
%% The following macros are setup in JFM.cls for sans-serif fonts in text
%% and math.  If you use these macros in your article, the required fonts
%% will be substitued when you article is typeset by the typesetter.
%%
%% \textsfi, \mathsfi   : sans-serif slanted
%% \textsfb, \mathsfb   : sans-serif bold
%% \textsfbi, \mathsfbi : sans-serif bold slanted (doesnt exist in CM fonts)
%%
%% For san-serif roman use \textsf and \mathsf as normal.
%%
%\newcommand\ssC{\mathsf{C}}    % for sans serif C
%\newcommand\sfsP{\mathsfi{P}}  % for sans serif sloping P
%\newcommand\slsQ{\mathsfbi{Q}} % for sans serif bold-sloping Q
%
%% Hat position
%\newcommand\hatp{\skew3\hat{p}}      % p with hat
%\newcommand\hatR{\skew3\hat{R}}      % R with hat
%\newcommand\hatRR{\skew3\hat{\hatR}} % R with 2 hats
%\newcommand\doubletildesigma{\skew2\tilde{\skew2\tilde{\Sigma}}}
%%       italic Sigma with double tilde
%
%% array strut to make delimiters come out right size both ends
%\newsavebox{\astrutbox}
%\sbox{\astrutbox}{\rule[-5pt]{0pt}{20pt}}
%\newcommand{\astrut}{\usebox{\astrutbox}}
%
%\newcommand\GaPQ{\ensuremath{G_a(P,Q)}}
%\newcommand\GsPQ{\ensuremath{G_s(P,Q)}}
%\newcommand\p{\ensuremath{\partial}}
%\newcommand\tti{\ensuremath{\rightarrow\infty}}
%\newcommand\kgd{\ensuremath{k\gamma d}}
%\newcommand\shalf{\ensuremath{{\scriptstyle\frac{1}{2}}}}
%\newcommand\sh{\ensuremath{^{\shalf}}}
%\newcommand\smh{\ensuremath{^{-\shalf}}}
%\newcommand\squart{\ensuremath{{\textstyle\frac{1}{4}}}}
%\newcommand\thalf{\ensuremath{{\textstyle\frac{1}{2}}}}
%\newcommand\Gat{\ensuremath{\widetilde{G_a}}}
%\newcommand\ttz{\ensuremath{\rightarrow 0}}
%\newcommand\ndq{\ensuremath{\frac{\mbox{$\partial$}}{\mbox{$\partial$} n_q}}}
%\newcommand\sumjm{\ensuremath{\sum_{j=1}^{M}}}
%\newcommand\pvi{\ensuremath{\int_0^{\infty}%
%  \mskip \ifCUPmtlplainloaded -30mu\else -33mu\fi -\quad}}
%
%\newcommand\etal{\mbox{\textit{et al.}}}
%\newcommand\etc{etc.\ }
%\newcommand\eg{e.g.\ }
%
%\newcommand\VA{\ensuremath{{\cal{V}}}}
%\newcommand\AV{\ensuremath{{\cal{A}}}}
%
%
%
%% Figures..
%\pgfplotsset{filter discard warning=false} % Emp?che l'affichage des warning lorsqu'un NaN est pr?sent dans le csv
\usepackage{textcomp}
